% Complete proofs for the Fraisse1954 presentation and its modern bridges.

\begin{annexitem}{annex:p3.fraisse54.basic}{thm:p03-fraisse54-basic}{Restrictions, chains, and the two embedding comparisons}
\begin{proof}
Identity maps and compositions prove reflexivity and transitivity of
\(\ll\). For transitivity of \(\prec\), take a finite restriction of
\(A\), embed it into \(B\), and then embed its finite image into \(C\).
Transport by isomorphisms preserves both constructions. An embedding
restricts to every finite subset, proving the implication between them.

For a chain, every finite tuple of elements in the union lies in one chain
member: finitely many chosen members have a largest one under inclusion.
Define the tuple's value there. A larger member has the same value by the
extension assumption, so the definition is independent of the member chosen.
It is the unique relation restricting to every chain member. The claim about
at most \(n\) elements follows because an \(n\)-tuple has at most \(n\)
distinct entries.

For a partial or linear order, the strict relation is the reflexive relation
with its diagonal deleted, and the reflexive relation is the strict relation
with the diagonal adjoined. Injections preserve and reflect equality, so
preserving and reflecting either order relation is equivalent to preserving
and reflecting the other. Thus the induced embeddings and automorphisms
agree under this translation.

For failure of the converse comparison, \((\mathbb N,<)\) and its reverse
have the same finite linear orders as restrictions but neither embeds in the
other: an embedding would yield an infinite descending sequence in a
well-order. For mutual embeddings without isomorphism, let
\(L=\sum_{k\in\mathbb N}\omega^*\), with the summands ordered by \(k\),
and let \(L'=1+L\). The order \(L\) has no least element, whereas \(L'\)
does. Inclusion embeds \(L\) into \(L'\). To embed \(L'\) into \(L\),
put its new first point in the first summand and its successive
\(\omega^*\)-summands in the second, third, and later summands. These are
the two distinct phenomena exhibited in the source's \S1.3.

The finite members of \(\Gamma_R\) are exactly \(\gamma_R\), so its
finite part determines the full class by the definition of \(\prec\).
If \(R\in\Gamma_S\), all its finite restrictions lie in \(\gamma_S\).
If every finite member of \(\Gamma_S\) also embeds into \(R\), the
reverse age inclusion holds and the two local classes coincide. Conversely
equality of local classes gives that finite embedding condition directly.
\end{proof}
\end{annexitem}

\begin{annexitem}{annex:p3.fraisse54.age-count}{thm:p03-fraisse54-age-count}{Counting the single-relation ages}
\begin{proof}
On a fixed \(p\)-element base an \(n\)-ary relation has
\(2^{p^n}\) possible value tables, so there are finitely many isomorphism
types of that size and countably many finite types altogether. Every age is
a subset of this countable collection, giving the upper bound
\(2^{\aleph_0}\).

A unary relation is a partition into positive and negative elements. Its
age is determined by the two capacities for finite positive and negative
subsets: each capacity is a natural number or \(\infty\), where
\(\infty\) means that all finite sizes occur. Conversely these capacities
determine which pairs of finite sizes occur, since disjoint positive and
negative selections can be made together. There are countably many capacity
pairs and infinitely many distinct finite positive capacities. Thus the unary
count is exactly \(\aleph_0\).

For arity two, take for each positive integer \(r\) a directed cycle
\(C_r\) with edges only to the next point; \(C_1\) has its single loop.
For \(S\subseteq\mathbb N_{>0}\), let \(A_S\) be the disjoint union of
the cycles with lengths in \(S\), with no cross-component edges. The
induced directed cycle \(C_r\) occurs in \(A_S\) exactly when \(r\in S\).
Indeed, its successor edges force its image to remain in one component,
and a nonempty proper restriction of a cycle does not contain that complete
cycle. Hence different \(S\) give different ages. For \(n>2\), replace
a binary relation \(A\) by the relation
\(A^*(x_1,\ldots,x_n)=A(x_1,x_2)\). A purported isomorphism of nonempty
restrictions of \(A^*,B^*\) preserves the binary relations by filling the
remaining coordinates with the second entry. Thus the distinct binary ages
remain distinct. Finally \(\gamma_R=\gamma_S\) exactly when
\(\Gamma_R=\Gamma_S\), since the latter classes are defined by those
finite members. This proves the source's Theorem I.
\end{proof}
\end{annexitem}

\begin{annexitem}{annex:p3.fraisse54.chain-hosts}{thm:p03-fraisse54-chain-hosts}{Successor and limit stages for a family of joint hosts}
\begin{proof}
For a finite or countable list \(A_i\), start with its first member and
at stage \(i+1\) choose a joint host of the preceding stage and \(A_i\).
Transport the host so the preceding stage is literally a restriction.
This gives the chain; each \(A_i\) embeds into a stage. The assertion
does not require the union itself to be in \(\mathcal K\).

For an arbitrary set-indexed list, Choice well-orders its index set.
Proceed at successors as before, and at a limit take the union of the
preceding chain before handling the next listed structure. Chain closure
puts that union in \(\mathcal K\), so joint embedding can be used again.
Isomorphism closure justifies every transport that fixes the preceding
base. If the list is empty use the empty chain.

For a definable class \(\mathcal K\), the witness selections in this
recursion can be confined to a set. To see this, let \(\lambda\) be its
set-sized length and choose an initial rank bound containing the given
family. From a rank bound \(\rho\), Collection gives a higher rank
containing joint-host witnesses for every pair in the set
\(\mathcal K\cap V_\rho\); include their transports and the unions of
all \(\lambda\)-indexed chains from that rank. At a limit bound the
union of the earlier rank stages and repeat. Transfinite recursion for
the unique least adequate rank bounds yields a final set containing all
required witnesses. Well-order that set once, and select its least
permitted witness at each actual construction step. This uses ordinary
Choice and Collection, not a global well-order of the universe.
\end{proof}
\end{annexitem}

\begin{annexitem}{annex:p3.fraisse54.age-realization}{thm:p03-fraisse54-age-realization}{Hereditary and joint-embedding characterization of ages}
\begin{proof}
An age is hereditary. For finite restrictions of a common relation, the
restriction to the union of their images supplies a joint embedding.
Conversely enumerate representatives \(A_0,A_1,\ldots\) of all the finite
types in \(\mathcal K\); a finite list can be repeated. Starting with a
member \(B_0\), use joint embedding to construct a finite member containing
a copy of \(B_i\) and \(A_i\). Rename it so the copy of \(B_i\) is
literally \(B_i\), and call it \(B_{i+1}\). The union \(R\) is finite
or countable. Every finite subset of the union lies in a stage, so heredity
gives \(\gamma_R\subseteq\mathcal K\). Every enumerated type embeds in
a stage, giving the reverse inclusion. No extension requirement over a
specified embedding has been used.

For a nonempty chain of ages, the union is hereditary and has joint
embedding by placing two of its members in the larger of their two ages.
The intersection is hereditary and contains the empty restriction. Let
\(A,B\) belong to the intersection. A joint embedding may always be
restricted to the union of the two images, of size at most
\(\lvert A\rvert+\lvert B\rvert\). There are finitely many possible
isomorphism types of such a joint host. The sets of permissible host types
for the individual ages form a nonempty decreasing chain of nonempty subsets
of that finite set. Their intersection is nonempty: choose a member of the
chain with minimum cardinality, after which further shrinking is impossible.
This supplies joint embedding in the intersection, so the characterization
already proved applies. This also handles an intersection having only the
empty relation.
\end{proof}
\end{annexitem}

\begin{annexitem}{annex:p3.fraisse54.filter-calculus}{thm:p03-fraisse54-filter-calculus}{The projection-filter operations}
\begin{proof}
If two inverse projections disagree at one tuple, the corresponding positive
and negative sets of maps both belong to the filter, but their intersection
is empty. This contradicts properness. For an ultrafilter, those two sets are
complementary and exactly one belongs to it; choose that sign at each tuple.
This defines the inverse projection and proves existence. Refinement retains
every set of maps witnessing a tuple's value.

Let \(r:D^E\to D^H\) be restriction. It is surjective: a map on \(H\)
extends by assigning one fixed element of the nonempty set \(D\) elsewhere.
Precisely,
\[
 Q\in\mathcal F_H\quad\Longleftrightarrow\quad r^{-1}[Q]\in\mathcal F.
\]
Inverse images preserve intersections, complements and the full and empty
sets, proving the filter laws and the ultrafilter assertion. Each tuple in
\(H\) has the same sign-set before and after restriction, proving the
restriction identity for inverse projections. If \(F\subseteq H\) is
finite, the preimage of the set of maps injective on \(F\) is exactly
\(I_F\subseteq D^E\). This proves preservation of normality by restriction;
preservation by refinement is immediate from membership.

When \(D\) is infinite, every finite subset of \(E\) injects into \(D\),
and such a map extends to all of \(E\) by a fixed default value. Thus each
\(I_F\) is nonempty, and
\(I_{F\cup G}\subseteq I_F\cap I_G\). The sets \(I_F\) therefore
generate a proper normal filter. Its ultrafilter refinement exists under the
stated ultrafilter principle, in particular under Choice. Explicitly, order
the proper filters extending a given proper filter by inclusion. The union
of a nonempty chain is again a proper filter: two of its members occur
together at a chain stage, and no stage contains the empty set. Zorn's
lemma gives a maximal proper extension \(\mathcal U\). If neither
\(Y\) nor its complement belongs to \(\mathcal U\), every
\(U\in\mathcal U\) meets \(Y\); otherwise upward closure would put
the complement in \(\mathcal U\). Thus adjoining \(Y\) generates a
proper larger filter, a contradiction. Properness prevents both
complements from belonging. Hence this maximal extension is an ultrafilter
in the stated definition.
\end{proof}
\end{annexitem}

\begin{annexitem}{annex:p3.fraisse54.filter-equality}{thm:p03-fraisse54-filter-equality}{Equality classes of maps and relation transport}
\begin{proof}
Reflexivity and symmetry follow from equality. If the sets witnessing
\(a\sim_{\mathcal F}b\) and \(b\sim_{\mathcal F}c\) belong to the
filter, their intersection is contained in the set witnessing
\(a\sim_{\mathcal F}c\). This proves transitivity. The restriction formula
for filters gives the same equivalence on any restricted domain.

For distinct \(a,b\), normality puts the inequality set in the filter,
which excludes the complementary equality set. Conversely, an ultrafilter
with no distinct equivalent elements contains each such inequality set.
Intersecting the finitely many inequalities between pairs of a finite set
gives its injectivity set, proving normality.

For the transport assertion, fix \(a_1,\ldots,a_n\in H\). Intersect the
sets witnessing \(f(a_i)=f(g(a_i))\) with the two sets witnessing the
values of \(A(\bar a)\) and \(A(g\bar a)\) as inverse projections.
This finite intersection is nonempty. Any one map in it shows that the two
values are equal. Since \(g\) is a bijection onto its image, it is the
required isomorphism. This supplies all equality and transport steps of
the source's \S\S3.9--3.10.
\end{proof}
\end{annexitem}

\begin{annexitem}{annex:p3.fraisse54.local-projection}{thm:p03-fraisse54-local-projection}{Finite restrictions determine projection filters}
\begin{proof}
Suppose \(A=\mathcal F^{-1}(R)\), and fix finite \(H\subseteq E\).
There are only finitely many tuples in \(H^n\). Intersect their sign-sets
in \(\mathcal F\), and choose a map in the nonempty intersection. Its
restriction to \(H\) is a surjection onto its image whose inverse relation
is exactly \(A|H\). If the filter is normal, intersect also with \(I_H\);
the chosen map is then an embedding.

Conversely, for each finite \(H\subseteq E\), let \(U_H\) be the set of
all maps \(E\to D\) whose restriction to \(H\) gives the required
inverse relation. Such a set is nonempty by the hypothesized finite map and
extension by a default value. Moreover
\(U_{H\cup J}\subseteq U_H\cap U_J\). The \(U_H\) generate a proper
filter. For any tuple \(\bar a\), take its finite set of entries \(H\);
then \(U_H\) is contained in the sign-set of that tuple, so
\(A=\mathcal F^{-1}(R)\). In the finite-embedding case require the maps
in \(U_H\) to be injective on \(H\). The same argument works, and
\(U_H\subseteq I_H\) makes the resulting filter normal. Ultrafilter
refinement preserves the inverse projection and normality.
\end{proof}
\end{annexitem}

\begin{annexitem}{annex:p3.fraisse54.gamma}{thm:p03-fraisse54-gamma}{Arbitrary bases, local closure, and common hosts}
\begin{proof}
The class \(\Gamma_R\) is hereditary and isomorphism closed. It is closed
under chain unions because each finite subset of a union lies in a stage.
If the base of \(R\) is finite, no member of \(\Gamma_R\) has more
elements than \(R\), and each member itself embeds in \(R\). If its
base \(D\) is infinite, take a normal ultrafilter on maps \(E\to D\)
for any prescribed set \(E\). Its inverse projection belongs to
\(\Gamma_R\) by the preceding theorem. This proves the size alternative.

We prove joint embedding by the original projection-filter construction.
The case of empty \(R\) is immediate. Otherwise put two members \(A,B\)
on disjoint bases \(F,G\). Choose normal filters on \(D^F,D^G\) giving
\(A,B\) as inverse projections. On \(D^{F\cup G}\), the preimages of
members of these two filters have the finite-intersection property:
independently chosen maps on the disjoint bases combine to one map. Extend
the generated filter to an ultrafilter \(\mathcal U\), and let
\(P=\mathcal U^{-1}(R)\). Its restrictions are \(A,B\), and
\(\sim_{\mathcal U}\) identifies no two distinct points of either
\(F\) or \(G\). Thus each equivalence class has at most two elements,
one in each base.

Let \(L\subseteq G\) be the points equivalent to a point of \(F\), and
put \(H=F\cup(G\setminus L)\). The restriction of \(\mathcal U\) to
\(H\) identifies no distinct points, so it is normal. Consequently
\(C=P|H\) belongs to \(\Gamma_R\). Inclusion embeds \(A\) into \(C\).
Send each point of \(L\) to its unique equivalent point in \(F\), and
fix \(G\setminus L\). This injective map is an isomorphism of \(B\)
onto its image by the filter-equality transport lemma. It embeds \(B\)
into \(C\), proving \(\mathrm C_2\). Notice that identifying points on
opposite sides was essential; this argument does not prove strong
amalgamation.

Now let a nonempty \(\mathcal K\) satisfy the three conditions.
Its finite members are hereditary and have joint embedding: restrict a
joint host to the two finite images. They therefore form an age
\(\gamma_S\). It remains to prove finite character. If a relation \(A\)
has all finite restrictions in \(\mathcal K\) but is not a member,
choose such a counterexample of least infinite cardinality \(\kappa\).
Well-order its base with order type the initial ordinal \(\kappa\).
Every proper initial segment has smaller cardinality, and its finite
restrictions belong to \(\mathcal K\), so minimality puts that segment
in \(\mathcal K\). Their chain union is \(A\), contradicting
\(\mathrm C_3\). Thus finite character holds, and it gives
\(\mathcal K=\Gamma_S\) in both directions.

Finally well-order a set-indexed family \((A_i)_{i<\lambda}\) in
\(\Gamma_R\). At each successor stage take a joint host of the preceding
stage and \(A_i\), renamed to extend the preceding stage literally.
At a limit take the chain union. Heredity allows each successor host to
be restricted to the two images, so all choices can be made from a set of
relations of bounded cardinality; no choice from a proper class is required.
The final union belongs to \(\Gamma_R\) and contains an image of every
\(A_i\). For the empty family choose any member. This proves Theorem IV,
including the arbitrary-index, successor and limit cases.
\end{proof}
\end{annexitem}

\begin{annexitem}{annex:p3.fraisse54.extension}{thm:p03-fraisse54-extension}{One-point extension, universality, and back-and-forth}
\begin{proof}
One-point extension implies the general finite extension property by adjoining
the finitely many new elements one at a time. Conversely a one-point problem
is a finite problem. If \(R\) is homogeneous, choose an embedding
\(j:B\to R\). For a specified \(e:A\to R\), the finite isomorphism
\(e(j|A)^{-1}\) extends to an automorphism \(g\) of \(R\). Then
\(gj:B\to R\) extends \(e\), proving \(\mathcal H\).

If \(A\prec R\) is countable, enumerate its base and let \(A_m\)
be its first \(m\) elements, starting at the empty restriction. Repeatedly
apply \(\mathcal H\) to extend an embedding of \(A_m\) to one of
\(A_{m+1}\). The union is injective and preserves every relation tuple
because that finite tuple lies in one stage. It embeds \(A\) in \(R\).
The same construction terminates for a finite base. No enumeration of the
target was needed.

For uniqueness, enumerate the bases of both countable structures and start
with the specified finite partial isomorphism \(f\). In a forth step,
take the first source point outside its domain. The finite restriction
formed by adjoining it belongs to the common age, so \(\mathcal H\)
in the target extends \(f\). In a back step apply the same argument to
the inverse partial isomorphism and the first target point outside its
domain. Invert again to obtain the next forward partial isomorphism.
If a chosen point is already included, keep the map unchanged. The union
is a bijection and preserves and reflects all relation tuples. It extends
the prescribed initial map. Finite structures of the same age have equal
size and the same argument ends with a full isomorphism. Taking the two
structures equal proves \(\mathcal H\Rightarrow\) homogeneity in the
finite or countable case. Together with the first paragraph this proves
all four assertions and the original Theorem V.
\end{proof}
\end{annexitem}

\begin{annexitem}{annex:p3.fraisse54.amalgamation}{thm:p03-fraisse54-amalgamation}{Amalgamation, its stronger form, and simultaneous finite extension}
\begin{proof}
Given \(\mathcal D\) and an embedding \(e:A'\to B\), with
\(A'\subseteq A\), amalgamate the inclusion into \(A\) with \(e\).
Rename the amalgam so the embedding of \(B\) is inclusion. The embedding
of \(A\) is the required extension, proving \(\mathcal D'\).
Conversely, identify the common domain of an amalgamation problem with
its image in \(A\), apply \(\mathcal D'\) to its embedding into
\(B\), and obtain the two embeddings into one extension of \(B\).
Renaming the structures outside the specified common domain prevents
accidental identifications of pre-existing names in these arguments.

For the stronger condition, first take copies whose bases intersect only
in the specified common part, then apply \(\overline{\mathcal D}\).
To see that the implication is strict, use the unary class of structures
with at most one positive element. It is \(\Gamma_R\) for a unary
relation \(R\) with one positive and infinitely many negative elements.
It has amalgamation: if both sides have a positive point outside the
common part, identify these two points; all other new points remain
distinct. If a positive point is already in the common part, both
embeddings preserve it. If one or both sides have no positive point,
take the union with no additional identification. These constructions
preserve both signs and yield at most one positive point. But two disjoint
positive singletons have no literal common extension in the class, since
such an extension would have two positive elements. Thus strong
amalgamation fails.

For a finite batch over \(R_0\), solve the first extension problem using
\(\mathcal D'\), obtaining \(R_1\supseteq R_0\). Compose the next
specified base embedding with inclusion into \(R_1\), solve that problem
and obtain \(R_2\). Continue through the finite list. Previous witnesses
remain embeddings because every later stage is an induced extension.
The last finite stage solves the entire list.
\end{proof}
\end{annexitem}

\begin{annexitem}{annex:p3.fraisse54.original-construction}{thm:p03-fraisse54-original-construction}{The original finite-size schedule and homogeneous representative}
\begin{proof}
Let \(\mathcal A\) be the finite members of \(\mathcal K\). They have
heredity and joint embedding. If a homogeneous relation has this age,
amalgamation follows by putting both finite structures inside it and
extending the isomorphism between the two images of the common domain
to an automorphism. Restrict to the union of the aligned images. This
proves necessity.

For sufficiency set \(R_0\) equal to the empty restriction. At stage
\(p\geq1\), choose one representative of each type in \(\mathcal A\)
of size at most \(p\). There are finitely many. For every such
\(B\), every restriction \(A\subseteq B\) of size at most \(p-1\),
and every embedding \(e:A\to R_{p-1}\), list the problem of extending
\(e\) to \(B\). This is a finite list: the representatives, their
subsets and the maps between finite bases all have finite possibilities.
Solve the list by the finite-batch lemma and call the result \(R_p\).
Including the empty restriction ensures that every listed \(B\) embeds
in \(R_p\). Equivalently, the source initializes the first stage by
jointly embedding the finitely many one-element types and then treats
the nonempty base embeddings. The empty-restriction convention makes
the degenerate empty age explicit.

Let \(R_0^*=\bigcup_pR_p\). Every finite subset lies in a stage, so
heredity gives \(\gamma_{R_0^*}\subseteq\mathcal A\). Every finite
\(B\in\mathcal A\) is listed once \(p\geq\lvert B\rvert\), proving
the reverse inclusion. For a specified extension problem
\(e:A\to R_0^*\), \(A\subseteq B\in\mathcal A\), choose \(m\)
with \(e[A]\subseteq R_m\), and choose
\[
 p\geq\max\{m+1,\lvert B\rvert,\lvert A\rvert+1\}.
\]
After replacing \(B\) by its enumerated isomorphic copy, this very
embedding problem occurs at stage \(p\). Its solution in \(R_p\)
solves the original problem in the union. Thus \(R_0^*\) has
\(\mathcal H\), and countability gives homogeneity.

The finite age equality implies \(\Gamma_{R_0^*}=\mathcal K\).
Back-and-forth gives uniqueness among finite or countable homogeneous
representatives. Every other homogeneous representative satisfies
\(\mathcal H\), so the countable universal embedding argument embeds
\(R_0^*\) into it. Applying the same argument with target \(R_0^*\)
embeds every finite or countable member of \(\mathcal K\). The proof
does not assert a countable host for an arbitrary family of arbitrary
uncountable members.
\end{proof}
\end{annexitem}

\begin{annexitem}{annex:p3.fraisse54.examples}{thm:p03-fraisse54-examples}{Relation, linear-order, partial-order, and ramified-order amalgams}
\begin{proof}
The property of being an arbitrary relation, a linear order, a partial
order, or a ramified partial order is determined by finite restrictions.
Their finite classes are hereditary and have joint embedding. For
relations choose arbitrary cross-tuple values on a disjoint union; for
partial and ramified orders take the disjoint union with no comparisons
between the sides; for linear orders place one side before the other.
The age-realization theorem and finite character make each full class
a \(\Gamma\)-class.

For arbitrary relations agreeing on a common base, assign the old values
on tuples entirely on one side and assign, for instance, minus on every
remaining mixed tuple. Agreement on the overlap makes this well-defined.
For linear orders, arrange the common finite base in increasing order.
In each gap, concatenate the new points from the first side, in their old
order, with the new points from the second side, in their old order.
Do the same before the first common point and after the last. If the
overlap is empty use a sum of the two orders. This is a literal common
extension.

For partial orders \(A,B\) on \(F,G\), with common induced order on
\(H=F\cap G\), take the transitive closure of
\(\leq_A\cup\leq_B\) on \(F\cup G\). It is reflexive and transitive.
Every finite path with endpoints in \(F\) can have each excursion through
\(G\) replaced by one comparison in \(F\): the excursion's endpoints
lie in \(H\), and the two orders agree there. Thus the resulting order
adds no comparison within \(F\), and symmetrically none within \(G\).
A cross comparison from \(a\in F\setminus H\) to
\(b\in G\setminus H\) holds exactly when some \(h\in H\) satisfies
\(a\leq_Ah\leq_Bb\), by the same path compression.
If also \(b\leq_Ca\), choose \(k\in H\) with
\(b\leq_Bk\leq_Aa\). We would have
\[
 a\leq_Ah,\qquad h\leq_Bb\leq_Bk,\qquad k\leq_Aa.
\]
Agreement on \(H\) gives \(h\leq_Ak\), so antisymmetry in \(A\)
forces \(a=h=k\), contradicting \(a\notin H\). This is the full
cross-side cycle argument represented in the source's Figure 1.
Hence the transitive closure is a partial order and a literal amalgam.

For ramified orders use \cref{fig:p03-fraisse54-ramified}. In \(A\), put
\(d<a,d<b\) and leave \(c\) incomparable with \(d\). In \(B\), put
\(e<b,e<c\) and leave \(a\) incomparable with \(e\). The common
\(a,b,c\) form an antichain; both four-point orders are ramified.
In any amalgam fixing that common part, the images of \(d,e\) are both
predecessors of \(b\). Ramification makes them comparable, including
possible equality. If \(d'\leq e'\), then \(d'<c\), contradicting
reflection of incomparability by the embedding of \(A\). If
\(e'\leq d'\), then \(e'<a\), contradicting the embedding of \(B\).
Thus even ordinary amalgamation fails. The positive homogeneous and
universal consequences, and the negative consequence for the full
ramified age, follow from the construction and necessity of amalgamation.
\end{proof}
\end{annexitem}

\begin{annexitem}{annex:p3.fraisse54.orders}{thm:p03-fraisse54-orders}{Homogeneous linear orders and the uncountable examples}
\begin{proof}
An increasing bijection between two finite subsets extends to an
automorphism exactly when corresponding gaps can be matched: the initial
segments, the intervals between consecutive selected points, and the final
segments must be isomorphic. The order sum of the chosen gap isomorphisms
and the finite bijection then gives the automorphism. Homogeneity itself
forces these isomorphisms by extending the one-point and two-point maps.
This proves the gap criterion.

Suppose an order of at least two elements satisfies \(\mathcal H\).
An endpoint would obstruct extending a one-point embedding to the
two-point order on its missing side, so endpoints do not exist. If
\(a<b\) were adjacent, choose a third point using absence of endpoints,
say \(c>b\). Map the two-point restriction on \(a,c\) to the adjacent
pair \(a,b\). It cannot extend to the three-point restriction containing
the intervening original \(b\). Thus the order is dense. Conversely,
to extend a finite embedding by one point, place the new image in the
appropriate initial gap, final gap or gap between consecutive images.
Absence of endpoints and density supply the image. Empty and one-point
orders separately satisfy homogeneity.

All nontrivial dense orders without endpoints have the finite linear
orders as their age. The countable extension/back-and-forth result therefore
identifies every countable such order with \(\mathbb Q\).

Here are the three uncountable homogeneous examples in the source's
\S6.3, with the order sums specified. Let
\[
 L_Q=\sum_{\alpha<\omega_1}\mathbb Q,
 \qquad
 L_R=\mathbb R+\sum_{\alpha<\omega_1}(1+\mathbb R),
\]
where summands are ordered by their index, and use \(\mathbb R\) itself
as the third example. Every segment below a point and every open interval
of \(L_Q\) is a countable dense order without endpoints, hence has type
\(\mathbb Q\). Each final segment has type \(L_Q\): the remaining part
of the current rational summand is rational, and the remaining index tail
has order type \(\omega_1\). The gap criterion proves homogeneity.

Every segment below a point and bounded open interval of \(L_R\) is a
separable Dedekind-complete dense order without endpoints. It meets only
countably many summands; countable dense subsets of those summands give
separability. Suprema within a real summand exist there, and a cut reaching
the end of an initial collection of whole summands is filled by the first
point of the next \(1+\mathbb R\) summand. These observations also prove
Dedekind completeness for subsets bounded inside an open interval.
Such an order \(J\) is isomorphic to \(\mathbb R\), as follows. Choose a
countable dense suborder \(D\), and use rational back-and-forth to obtain
\(h:D\to\mathbb Q\). Define
\[
 t(x)=\sup\{h(d):d\in D,\ d<x\}.
\]
The set is nonempty and bounded by choosing members of \(D\) on either
side of \(x\). If \(x<y\), choose \(a,b\in D\) with
\(x<a<b<y\); then \(t(x)\leq h(a)<h(b)\leq t(y)\), proving strict
increase. For \(r\in\mathbb R\), the set
\(A_r=\{d\in D:h(d)<r\}\) is nonempty and bounded in \(J\); let
\(x=\sup_JA_r\). If \(d<x\), some \(a\in A_r\) satisfies \(d<a\),
since otherwise \(d\) would be a smaller upper bound; hence \(h(d)<r\).
Conversely, for rational \(q<r\), choose a rational \(q'\) strictly
between them. Then \(h^{-1}(q)<h^{-1}(q')\leq x\), so \(t(x)\geq q\).
It follows that \(t(x)=r\). Thus \(t\) is an increasing bijection.
Every final segment of \(L_R\) has type \(L_R\), since its first real
tail has type \(\mathbb R\) and the remaining index tail is
\(\omega_1\). Again the gap criterion applies. It also applies directly
to \(\mathbb R\). The orders are distinct: \(L_Q\) has countable
initial segments, whereas the other two do not; \(\mathbb R\) has
countable cofinality, whereas \(L_R\) has cofinality \(\omega_1\).

Finally \(\mathbb R\setminus\{0\}\) is dense without endpoints and
therefore has \(\mathcal H\). For a negative point its initial segment
is Dedekind complete. For a positive point the corresponding segment has
the unfilled cut at zero: its negative elements form a bounded set without
a least upper bound. The initial segments are not isomorphic, so the gap
criterion refutes homogeneity.
\end{proof}
\end{annexitem}

\begin{annexitem}{annex:p3.fraisse54.countable-language}{thm:p03-fraisse54-countable-language}{A fair schedule for countably many relational types}
\begin{proof}
The necessity of heredity and countably many types follows by listing the
finite subsets of a countable structure. Unbounded sizes follow from its
infinite base, joint embedding from union of finite images, and amalgamation
from an automorphism aligning the two embeddings of a common finite part.

Conversely enumerate representatives \(B_0,B_1,\ldots\) of
\(\mathcal K\). There is a unique empty structure in \(\mathcal K\):
if nullary relation symbols are present, joint embedding forces all members
to give them the same values. Start with this empty structure \(M_0\).
At stage \(s\), form the finite list of all triples
\[
 (i,A,e),\qquad i\leq s,\quad A\subseteq B_i,
 \quad e:A\longrightarrow M_s\text{ an embedding}.
\]
There are finitely many indices, finitely many subsets of each finite
\(B_i\), and finitely many injections into the finite \(M_s\).
Successively amalgamate \(B_i\) and the current extension over \(A\),
using the specified \(e\) followed by inclusion. Rename the new amalgam
to fix that current extension. After the finite list, call the result
\(M_{s+1}\). Heredity permits restricting each amalgam to the two images,
so all stages are finite members of \(\mathcal K\).

For \(M=\bigcup_sM_s\), every finite restriction lies in a stage and
belongs to \(\mathcal K\). The empty-base problem for \(B_i\) is
included at every stage \(s\geq i\), so every representative embeds in
the union. Thus \(\operatorname{Age}(M)=\mathcal K\), and unbounded
sizes make \(M\) countably infinite.

Given a finite extension problem \(e:A\to M\), \(A\subseteq B\in
\mathcal K\), take an isomorphism \(h:B_i\to B\) and a stage \(N\)
containing \(e[A]\). For \(s\geq\max\{i,N\}\), the triple
\((i,h^{-1}[A],eh|h^{-1}[A])\) occurs in that stage's finite list.
Its solution, composed with \(h^{-1}\), extends the original \(e\).
This proves the full finite extension property. The earlier forth and
back-and-forth proofs use only finite induced structures, so they apply
unchanged to all the relation symbols at once. They give homogeneity,
uniqueness and countable universality. The schedule is an existence
construction, not a claim that membership in \(\mathcal K\) or the
embedding test is computable.
\end{proof}
\end{annexitem}

\begin{annexitem}{annex:p3.fraisse54.categoricity-boundary}{thm:p03-fraisse54-categoricity-boundary}{Why finiteness of the signature matters}
\begin{proof}
Fix a tuple length \(m\). In a finite relational signature, only finitely
many atomic formulas can be formed using its \(m\) variables, including
equalities and nullary atoms. A complete assignment of signs to these
atoms is a finite quantifier-free diagram. Two realized tuples with the
same diagram give an isomorphism between their induced finite substructures;
ultrahomogeneity extends it to an automorphism. Thus any formula on these
tuples is equivalent in the structure to a finite disjunction of their
realized atomic diagrams. The universal closure of this equivalence belongs
to its complete theory. This proves quantifier elimination, including the
zero-variable case using truth and falsity.

For categoricity, express each finite extension problem by the finite
diagrams of its domain and codomain. The sentence saying that every
realization of the domain diagram extends to the codomain diagram is true
in the homogeneous structure, hence in its complete theory. In any other
countable model of that theory, existential finite diagrams give exactly
the same age, and the extension sentences give \(\mathcal H\).
Back-and-forth supplies an isomorphism.

For the failure in a countable signature, take unary symbols
\(P_0,P_1,\ldots\). Let \(F\) be the disjoint union of countably many
countably infinite sets, with \(P_i\) naming its \(i\)-th set. Its age
is the class of finite sets whose elements each have exactly one of these
colors. There are countably many finite types. Heredity, unbounded sizes,
joint embedding and free amalgamation all hold. A finite partial
isomorphism preserves each color and extends to a permutation of every
infinite color class, so \(F\) is ultrahomogeneous.

Let \(F'\) add one uncolored point. For any finite sublanguage, the named
color classes in \(F,F'\) are countably infinite, and their complementary
sets are also countably infinite. Matching these finitely many named classes
and their complements gives an isomorphism of the reducts. Every
first-order sentence uses only finitely many symbols, so \(F,F'\) satisfy
the same sentences. But they are not isomorphic in the full language:
every point of \(F\) has a color, and the added point of \(F'\) has none.
Thus the complete theory of the countable Fra\"iss\'e limit \(F\) is not
\(\aleph_0\)-categorical.
\end{proof}
\end{annexitem}

\begin{annexitem}{annex:p3.fraisse54.graph}{thm:p03-fraisse54-graph}{The graph amalgam and every extension instance}
\begin{proof}
An induced subgraph of a finite simple graph is finite and simple. Disjoint
union supplies joint embedding. For an amalgamation problem, take copies
whose intersection is exactly the common induced subgraph. On their union
keep all old edges and add no edges between the two new parts. Agreement
on the common graph makes this an irreflexive symmetric relation, and both
inclusions preserve edges and nonedges. There are finitely many graph types
of each size, and sizes are unbounded. The limit theorem applies.

For finite disjoint \(U,V\) in the limit, take the induced graph on
\(U\cup V\), add one new vertex with edges exactly to \(U\), and
extend its inclusion by \(\mathcal H\). Injectivity ensures that the
new image is outside \(U\cup V\). Conversely the displayed property
extends an embedding by one vertex: place its already embedded neighbors
in \(U\) and its already embedded nonneighbors in \(V\). Every finite
extension follows by iteration. It also embeds every finite graph, and
the forth construction embeds every countable graph as an induced subgraph.

In the schema of \cref{def:p03-fraisse54-graph-axioms}, the new vertex is
adjacent to the first \(r\) and nonadjacent to the last \(s\) of any
\(r+s\) distinct vertices. These extension axioms imply
infinitude, since a fresh vertex can be required outside any finite list.
They give \(\mathcal H\) and the full finite graph age, so all countable
models are isomorphic.

For completeness without restricting to countable models, use finite
partial isomorphisms between any two models of the axioms. Atomic formulas
agree. Boolean combinations preserve agreement. For an existential formula,
an old witness already in the domain transfers along the map; a new witness
has a finite adjacency/nonadjacency pattern to the domain, which the target
extension axiom realizes outside the current range. Apply induction to the
subformula. The reverse direction uses the symmetric construction. Thus
all formulas agree on related tuples; starting with the empty map proves
that the two structures satisfy the same sentences. The extension theory
is satisfiable by the constructed limit and is complete, and its unique
countable model proves \(\aleph_0\)-categoricity.
\end{proof}
\end{annexitem}

\begin{annexitem}{annex:p3.fraisse54.logical-joint}{thm:p03-fraisse54-logical-joint}{A common elementary host by diagrams}
\begin{proof}
For each model \(M_i\), add a separately tagged constant for each of its
elements and take its full elementary diagram. Let \(\Sigma\) be the
union of these diagrams with \(T\). Consider a finite subset. For each
of its finitely many contributing models, replace its constants by distinct
variables to obtain a finite conjunction \(\delta_i(\bar x_i)\), true
of a tuple in \(M_i\). Completeness and satisfiability of \(T\) imply
\(T\models\exists\bar x_i\,\delta_i\). Hence a model of \(T\)
realizes all these finitely many existential requirements. The constant
sets are separately tagged, so interpreting their finite lists with those
witnesses satisfies the chosen finite subset of \(\Sigma\).

First-order compactness, \cref{thm:fol-compactness}, gives a model
of \(\Sigma\). The interpretations of the constants produce elementary
embeddings of all \(M_i\) into its reduct: inequalities in each diagram
give injectivity and all formula instances give elementarity. Distinct
models' images may overlap. With an empty family any model of \(T\)
is a host. This proves the claimed joint embedding and the stated
elementary strengthening of the source's logical application.
\end{proof}
\end{annexitem}
